\preludeandfugue{Disassociation}
\contributor{Dana Archer}

% Biography

% Dana is a fiction writer and film critic who is currently seeking to
% publish her debut novel. She is an avid reader and has written for
% many years in a number of formats from fan-fiction to short stories.

% This is all my own work and I hold all rights to it. It is written
% from experience and is not slanderous or misleading.

% Words chosen - Tense, Illusory, confusion, Fractured sense of reality, Disassociation 

\prelude

Disassociation is about as hard to describe as it is to spell. It is a
feeling of being somewhere physically but mentally somewhere entirely
different. Although I eventually realised that this feeling has a
name, it still perplexes me. Maladaptive daydreaming is an
increasingly common yet massively under-researched psychological
disorder where a person daydreams excessively and sometimes to a point
where reality becomes numb by comparison. Paracosm is the term used
for this other world, one that is crafted sometimes with aspects
borrowed from media and fandoms. Unfortunately a world crafted over
many years can become more a prison than an escape as it numbs the
constraints of reality with a place unburdened by logic.  It's
interesting that music is the prompt for this anthology as often it is
a key catalyst for the daydreams. It is impossible to ever fully
replicate the experience of being engulfed inside your own mind to the
point that it invokes a dissolution of reality.
% However, with this piece I aim to try and convey the feeling of
% being trapped inside a world you’ve created, a euphoria that has
% taken many years to build only to become a tomb. With it comes
% memories and moments from your life that you had spent so many years
% running from. Excessive daydreaming at its core is an escape, a
% reaction to several external factors such as Trauma or loneliness.
What really is this feeling of being trapped inside a
world you've created, a euphoria that has taken many years to build
only to become a tomb? And what memories come back? What moments from
a life that you had spent so many years running from? Excessive
daydreaming at its core is an escape, a reaction to several external
factors such as trauma or loneliness.

\fugue 

I didn't notice it at first. This was odd really, as it was just as
% perhaps the most
present then as any time later on. [This paragraph made no sense to me in the context of what comes later so I have rewritten it.]
% \marginpar{prevalent?} it had ever been.
To start with, it was just a constant trickle, a gentle drop gliding down the
glass hidden behind a curtain. Of course,
% as much as it was always there,
ignorance of its presence yields 
innocence, comfort yields longing. I struggle
to remember there ever being a time when these feelings weren't present, but there
must have been one. At some point there must have been silence too, yet I
fail to recall it.

I am sitting in a room full of my peers, yet only being able to hear
that dripping, tapping at the window. An infinitesimal sound is somehow
overpowering a room filled with voices. I once remember walking
through a crowded shopping centre during a bout of extreme heatwave;
people were rushing in all directions and I didn't notice them. They were packed like
sardines in that concrete catacomb of capitalist necessity and I was
oblivious of them. There was no dripping as the sun baked through the glass
ceiling, beads of sweat making puddles on the ground beneath us.
I was there physically, but my brain wandered, longing for the slightest
indication of rain. It traveled far that day in pursuit as I fought
through every step, desperate to stay present. But there doesn't have
to be rainfall to create condensation and the puddling heat proved an
unrelenting motivator.

I had built a world that my brain couldn't resist, an addictive tonic
to my own dismal reality.

As I aged, I became more aware. Escape doesn't serve me, nor will it
better my mental state. It might in the short term, perhaps, but not
in the long term. At first, I welcomed it, but it soon evolved into a
dangerous time-consuming crater. An addictive resort in which my mind
longed to revel in.

Around that time, the transparent teetering increased to a patter.
[I don't understand:  you've said it was just as prevant at first so how can it increase?
Since you seem sure it is increasing I changed what you wrote before.]
A tapping against the glass, concealable sometimes through music or
distraction. It never dissipated however, growing in strength,
demanding my attention. At times, when my attention was already
scattered between multiple needs, it screamed at me. Its steady
knocking evolving into a thud not unlike a demolition mission outside
my window. I wouldn't hide\=I couldn't. Although I happened to be
preoccupied, it waited silently for any brief interlude. A tempestuous
conflict with my brain and my existence. I had offered myself a
better alternative and my brain refused to part with it.

On a day I'm unable to pin-point, I allowed the droplets
in. Engulfing myself under their refreshing majesty. It was a moment
of sheer euphoria, a relief which I hoped would bring with it the
denouement to my torment.  As I stood beneath their dizzying waterfall
of oblivion, bathed in disappearing worries I little knew know what
the real world had turned its back on.

I've realized I was wrong to invite it in. What was once a patter had
become a wave. If I am lucky, I can hide away before the water
thickens. Today, however, I have found myself at the mercy of a tidal
wave. I used to run away, squeeze myself into a corner as the
rapturous aqua phenomenon washes away every aspect of life possible. I
learned quite early on that I couldn't out-run it and hiding from it
was illogical. So I simply stand, facing head on to that which has the power to
wash me into oblivion.

Oblivion wasn't so bad. [I don't understand.  By what you say in the last paragraph you haven't reached oblivion yet.]

As I allow the wave to take me, I swirl around in a
hurricane. Disorientated and lost, I try not to ponder or
guess. Navigating the numbing cold would be a waste of what little
energy I am permitted.  A limit of my own consequence. As the water
settles, I am somewhere in the middle of nowhere. This doesn't
frighten me as I look around for signs of life. [Ambiguous. Do you mean the search for signs of life
  prevents you being frightened, or that you are looking around for signs for life and simultaneously
are not frighted for a different reason?] Unafraid, I find
myself staring into the eye of a solitary hurricane resting above an
onyx underbelly. Confident that only I fell [fall?] privy to the wave, I look
down. The water splashes and shifts as I part it with my fingers. I
have no power here. Still, I squint aimlessly trying to see the
elegance embedded within the dark.

I decide to swim down.

Swimming forward would prove futile, there isn't anything to see above
the surface. [New paragraph here?] Beneath the darkening ocean is a collection of debris
and dust. Swimming further into the abyss, I find a glimmer of
light. A delicate speckle of illumination. Tiny glimmers upon the
rocks, my lungs tighten, it's becoming harder to breathe. The further
I go, the harder it is to move. I am sluggish, my muscles tighten. I
am becoming anchored to the deep.

Refusing to face defeat, I fight against the current. My body is
pleading to stop and yet I am not dissuaded. Continuing downwards. I
cannot die. It might tell me otherwise but I have grown skeptical of
the voice. Exquisite beauty has crystallized over years of
solitude. Untouched, unscathed, aged by its surroundings\=this is
new indeed.  Narrow at first, I am hunched into a crouch as I
barely make it through the small chasm.

I spare a thought for the world I have left behind. Sharp uneven walls
scrape and pull at my skin as I push past. Those who continue will
establish rules and beliefs for the many to follow. I only hope they
are kinder than before the wave. Dark algae coats the rocks, marking
my skin with lines and dusting. Perhaps I'll be a myth, a girl who
waited as the wave washed her away. A girl for whom the world wasn't
enough\=so she had to build one. A smell is developing, pebbles under
foot prove tricky to pass over.

`She never did fit anywhere,' neighbours will muse.  Perhaps they'll
write books, make documentaries\=media and culture often tries to
unpack what the public won't understand. Not that any of it matters to
me now. I became socially isolated due to fear\=a fear of what they'd
think, if I slipped up. This kept me safely in the grasp of my own
conscious, anchored to a world of my own crafting.

Turning a corner, I find a grandiose mausoleum. Embedded into the
walls are statues in varying states of decay. Faces unrecognizable,
some startling\=I squeeze my eyes shut but still can see them as clearly
as I had when they were open. I've gone too far.

When I allowed the rain in, I traded with it my control. Fear claws
under my skin, hooking into my muscles and diluting my veins. Freezing
to the spot, my muscles refuse to go any further.

I try to go back, I try to swim upwards. I am unable. My body has
hardened, I am rooted here only able to move forward but not
back. Taking a moment to try and catch my hastening breath, I lie
back. It doesn't help: nausea arises. Movement evolves around me, the
statues quake and quiver. There are no visible exits, just more of the
same. Finding my feet, I rush against the friction of the
water. Animated with jolts, the arms reach out to me, pulling me to
the centre. It is an orchestra of faces that I had forgotten, repressed in the
deepest depths of my mind, performing in unison. There is no sound, just
mechanical movements of puppets defying the rules of rigor mortis.

A prisoner in the best seat of the house, I clutch at my chest. This is a show I
created, one I had spent most of my life running from. Bio-luminescent
green charges up the walls, casting light upon the figures as they
each crawl from their stone prison. I see them more clearly, some are
completely skeletal\=the eyes are the only distinguishing feature. It
isn't those who frighten me, but the figures who are fresher. Flesh
clings where it is able; faces are almost entirely intact. These are the
ones who cause my heart to race and my stomach to tumble into
knots. They move as I remember, acting out slowly all that I've tried to
forget. Habitual and cold, they perform to an unwilling audience.
Unable to stand, the seat begins to sprout vines. As thick as they are
heavy, they bind me to the stone. A willing sacrifice, this is the price I
must pay for my own ignorance.

When you hide away in your own psyche for so long, it's inevitable that
it'll show you that which you were hiding from.

It crosses my mind that this time I just might die. A thought which
amuses me somewhat. They are only as strong as the power I afford them.

As soon as the thought flashes by, the figures slow to a halt. A glow
in the distance, for but a second. Pulling my way from the tightening
algae, I claw my way through the dense water.  It was real. I saw
it. I tell myself over and over as the grazes and tears sting from the
salt water.

I dare to look up to a flutter of daylight cracks through the kinetic
waves. Suddenly I am able to breathe: the air is clear, the water
thinner as I swim upwards. Drawing breath, I am shaky but full.

I open my eyes, the window is coated with condensation. The sky is a
welcoming dawn and the world continues. Turning away and drawing the
curtains, I climb into bed.

I know it won't be long before I go back there. 


